%%%%%%%%%%%%%%%%%%%%%%%%%%%%%%%%%%%%%%%%%%%%%%%%%%%%%%%%%%%%%%%%%%%%%%%%%%%%%%%
%  SÍLABO DE MICROECONOMETRÍA — Plan de estudio integrado
%
%  Textos base:
%    CT  — Cameron y Trivedi, Microeconometrics: Methods and Applications
%          (Cambridge, 2005)
%    CI  — Cunningham, Causal Inference: The Remix (2.ª ed., en línea,
%          https://mixtape.scunning.com)
%  Textos de apoyo:
%    MHE — Angrist y Pischke, Mostly Harmless Econometrics (Princeton, 2009)
%    MM  — Angrist y Pischke, Mastering 'Metrics (Princeton, 2015)
%    Wd  — Wooldridge, Econometric Analysis of Cross Section and Panel
%          Data (MIT Press, 2002)
%    Bt  — Baltagi, Econometric Analysis of Panel Data (Springer, 6.ª ed., 2021)
%    Ar  — Arellano, Panel Data Econometrics (Oxford, 2003)
%    Hu  — Huber, Causal Analysis (MIT Press, 2023)
%  Programas PUCP integrados: 1ECO25 (2025-1, J. M. Rentería; 2025-2,
%    L. García), 1ECO77 (2025-2, 2026-2), ECO315 (2022-0), ECO387 (2022-2),
%    ECO386 (2024-2) y la Unidad 3 de 2ECO01.
%
%  Versión revisada (octubre de 2026): ver «Cambios respecto de la versión anterior».
%  Compilar con:  pdflatex Microeconometria_revisado.tex   (dos veces, por el índice)
%%%%%%%%%%%%%%%%%%%%%%%%%%%%%%%%%%%%%%%%%%%%%%%%%%%%%%%%%%%%%%%%%%%%%%%%%%%%%%%

\documentclass[11pt,a4paper]{article}

\usepackage[T1]{fontenc}
\usepackage[utf8]{inputenc}
\usepackage[spanish,es-noquoting,es-tabla]{babel}
\usepackage{lmodern}
\usepackage[scaled=0.92]{helvet}        % títulos en sans serif, al estilo McGraw-Hill
\usepackage[margin=2.4cm,headheight=15pt]{geometry}
\usepackage[table]{xcolor}
\usepackage{enumitem}
\usepackage{titlesec}
\usepackage{titletoc}
\usepackage{tabularx}
\usepackage{booktabs}
\usepackage{longtable}
\usepackage{amsmath}
\usepackage{amssymb}
\usepackage{fancyhdr}
\usepackage{tikz}
\usepackage[skins]{tcolorbox}
\usepackage{needspace}

%------------------------------------------------------------------ colores ---
% Gama de verdes oliva
\definecolor{tonoOscuro}{RGB}{29,41,5}              % oscuro
\definecolor{tonoProfundo}{RGB}{54,77,10}           % profundo
\definecolor{tonoPrincipal}{RGB}{80,113,15}         % principal
\definecolor{tonoA}{RGB}{78,134,30}                 % variante cálida
\definecolor{tonoB}{RGB}{113,132,21}                % variante fría
\definecolor{tonoC}{RGB}{146,146,28}                % variante viva
\definecolor{tonoClaro}{RGB}{176,212,105}           % claro
\definecolor{tonoAcero}{RGB}{84,99,54}              % acero
\definecolor{tonoPizarra}{RGB}{95,106,73}           % pizarra
\definecolor{tonoGrafito}{RGB}{39,45,27}            % grafito (texto secundario)
\definecolor{tonoHielo}{RGB}{233,244,210}           % hielo
\definecolor{tonoPalido}{RGB}{245,250,235}          % pálido
\definecolor{tonoBruma}{RGB}{251,253,246}           % bruma

\colorlet{azulosc}{tonoOscuro}
\colorlet{grisclaro}{tonoPalido}
% Un tono distinto de la gama para cada libro
\definecolor{colorCT}{RGB}{94,92,18}
\definecolor{colorCI}{RGB}{129,137,26}
\definecolor{colorMHE}{RGB}{99,116,22}
\definecolor{colorMM}{RGB}{73,94,18}
\definecolor{colorWd}{RGB}{94,137,26}
\definecolor{colorBt}{RGB}{70,116,22}
\definecolor{colorAr}{RGB}{49,94,18}
\definecolor{colorHu}{RGB}{59,137,26}
\colorlet{colorPUCP}{tonoPizarra}

\usepackage[colorlinks=true,linkcolor=tonoOscuro,urlcolor=tonoPrincipal]{hyperref}

%------------------------------------------------------------------ estilos ---
\titleformat{\section}[block]
  {\sffamily\LARGE\bfseries\color{tonoOscuro}\raggedright}{\thesection.}{0.6em}
  {\textcolor{tonoPrincipal}{\rule[-0.1ex]{0.55em}{0.85em}}\hspace{0.45em}}
  [\vspace{-0.4ex}{\color{tonoPrincipal}\titlerule[1.6pt]}\vspace{1.2pt}{\color{tonoClaro}\titlerule[0.6pt]}]
\titlespacing*{\section}{0pt}{2.6ex plus 1ex minus .2ex}{1.6ex plus .2ex}

\titleformat{\subsection}
  {\sffamily\large\bfseries\color{tonoPrincipal}}{\thesubsection}{0.6em}
  {\textcolor{tonoClaro}{\rule[0.1ex]{0.45em}{0.45em}}\hspace{0.45em}}
\titlespacing*{\subsection}{0pt}{2.2ex plus .8ex minus .2ex}{1ex plus .2ex}

% Índice
\titlecontents{section}[0pt]
  {\addvspace{3pt}\sffamily\bfseries\color{tonoOscuro}}
  {}{\textcolor{tonoB}{\rule[0.15ex]{0.4em}{0.4em}}\hspace{0.6em}}
  {{\color{tonoClaro}\titlerule*[0.7em]{.}}\textcolor{tonoPrincipal}{\contentspage}}

% Cabecera y pie
\pagestyle{fancy}
\fancyhf{}
\lhead{\small\sffamily\bfseries\color{tonoOscuro}Sílabo de Microeconometría}
\rhead{\footnotesize\sffamily\color{tonoPrincipal}Cameron--Trivedi $\cdot$ Cunningham $\cdot$ Angrist--Pischke $\cdot$ Wooldridge $\cdot$ Baltagi}
\cfoot{{\setlength{\fboxsep}{3pt}\colorbox{tonoPrincipal}{\makebox[1.8em]{\small\sffamily\bfseries\color{white}\thepage}}}}
\renewcommand{\headrule}{%
  {\color{tonoPrincipal}\hrule height 1.1pt width\headwidth}\vskip 1.2pt
  {\color{tonoClaro}\hrule height 0.4pt width\headwidth}\vskip -2.7pt}

% Listas
\setlist[itemize,1]{label=\textcolor{tonoPrincipal}{\scriptsize$\blacksquare$}}
\setlist[enumerate,1]{label=\textcolor{tonoPrincipal}{\sffamily\bfseries\arabic*.}}

% Tablas: cabecera oscura y filas alternas en tono pálido
\arrayrulecolor{tonoPrincipal}
\setlength{\aboverulesep}{0pt}
\setlength{\belowrulesep}{0pt}
\renewcommand{\arraystretch}{1.3}
\newcommand{\cab}[1]{\textcolor{white}{\sffamily\bfseries #1}}
\newcommand{\filasalternas}{\rowcolors{2}{white}{tonoPalido}}

% Encabezados de párrafo en color
\newcommand{\destacado}[1]{{\sffamily\bfseries\color{tonoPrincipal}#1}}

% Franja decorativa con la gama
\newcommand{\franjas}[1][3pt]{%
  \noindent\begin{tikzpicture}
    \foreach \c [count=\i from 0] in {tonoOscuro,tonoProfundo,tonoPrincipal,tonoA,tonoB,tonoC,tonoClaro,tonoHielo}
      \fill[\c] (\i*\linewidth/8,0) rectangle ++(\linewidth/8,#1);
  \end{tikzpicture}\par}

%---------------------------------------------------------------- comandos ----
% Etiqueta de libro: recuadro de color con la sigla
\newcommand{\etiqueta}[2]{%
  {\setlength{\fboxsep}{1.4pt}\colorbox{#1}{\sffamily\bfseries\color{white}\strut#2}}}
% Referencias abreviadas a cada libro
\newcommand{\CT}[1]{\etiqueta{colorCT}{CT}\nobreak\hspace{0.25em}\textcolor{colorCT}{#1}}
\newcommand{\CI}[1]{\etiqueta{colorCI}{CI}\nobreak\hspace{0.25em}\textcolor{colorCI}{#1}}
\newcommand{\MHE}[1]{\etiqueta{colorMHE}{MHE}\nobreak\hspace{0.25em}\textcolor{colorMHE}{#1}}
\newcommand{\MM}[1]{\etiqueta{colorMM}{MM}\nobreak\hspace{0.25em}\textcolor{colorMM}{#1}}
\newcommand{\WD}[1]{\etiqueta{colorWd}{Wd}\nobreak\hspace{0.25em}\textcolor{colorWd}{#1}}
\newcommand{\BT}[1]{\etiqueta{colorBt}{Bt}\nobreak\hspace{0.25em}\textcolor{colorBt}{#1}}
\newcommand{\AR}[1]{\etiqueta{colorAr}{Ar}\nobreak\hspace{0.25em}\textcolor{colorAr}{#1}}
\newcommand{\HU}[1]{\etiqueta{colorHu}{Hu}\nobreak\hspace{0.25em}\textcolor{colorHu}{#1}}

% Un tema: número en recuadro de color + título
\newcommand{\tema}[1]{\temaaux#1\relax}
\def\temaaux#1 #2\relax{%
  \par\addvspace{1.1em}\noindent
  {\sffamily\bfseries\large\raggedright
   \hangindent=2.9em\hangafter=1
   \makebox[2.9em][l]{{\setlength{\fboxsep}{2.5pt}\colorbox{tonoPrincipal}{\color{white}#1}}}%
   {\color{tonoOscuro}#2}\par}%
  \nobreak\vspace{0.25em}\noindent{\color{tonoClaro}\rule{\linewidth}{0.6pt}}\par\nobreak\vspace{0.35em}}

% Línea de lecturas con marcador y etiqueta de color
\newcommand{\linealectura}[3]{%
  \par\nobreak\noindent\hspace{1.2em}%
  \begin{minipage}{0.92\textwidth}\small
    \textcolor{#1}{\rule[0.1ex]{0.45em}{0.45em}}\hspace{0.45em}%
    {\sffamily\bfseries\color{#1}#2}\enspace#3%
  \end{minipage}\par\vspace{0.3em}}
\newcommand{\lecturas}[1]{\linealectura{tonoPrincipal}{Lecturas:}{#1}}
\newcommand{\apoyo}[1]{\linealectura{tonoA}{Profundización:}{#1}}
\newcommand{\juegos}[1]{\linealectura{tonoB}{Teoría de juegos:}{#1}}
\newcommand{\supl}[1]{\linealectura{tonoC}{Suplemento:}{#1}}
\newtcolorbox{cajanota}{enhanced, sharp corners, boxrule=0pt, frame hidden,
  colback=tonoPalido, borderline west={2.5pt}{0pt}{tonoClaro},
  left=8pt, right=6pt, top=4pt, bottom=4pt, before skip=0pt, after skip=0pt,
  fontupper=\small\color{tonoGrafito}}
\newcommand{\nota}[1]{%
  \par\nobreak\noindent\hspace{1.2em}%
  \begin{minipage}{0.92\textwidth}\begin{cajanota}#1\end{cajanota}\end{minipage}\par\vspace{0.3em}}
\newcommand{\consultalib}[1]{\linealectura{tonoPizarra}{Consulta:}{{\color{tonoGrafito}#1}}}
% Artículos y textos tomados de los programas PUCP
\newcommand{\articulos}[1]{\linealectura{colorPUCP}{Artículos (PUCP):}{#1}}
% Artículos que no proceden de los programas PUCP (entran en las horas)
\newcommand{\articulosB}[1]{\linealectura{tonoB}{Artículos:}{#1}}
\newcommand{\consulta}[1]{\linealectura{colorPUCP}{Consulta (PUCP):}{{\color{tonoGrafito}#1}}}

% Encabezado de bloque: banda oscura, franja de la gama y recuadro de presentación
\newif\ifrecienparte
\newcommand{\bloque}[2]{%
  \ifrecienparte\global\recienpartefalse\else\par\needspace{0.4\textheight}\fi
  \phantomsection\addcontentsline{toc}{section}{#1}%
  \begin{tcolorbox}[enhanced, sharp corners, boxrule=0pt, frame hidden,
      colback=tonoOscuro, borderline west={7pt}{0pt}{tonoClaro},
      left=14pt, right=10pt, top=10pt, bottom=10pt, before skip=0pt, after skip=0pt]
    {\raggedright\sffamily\Large\bfseries\color{white}#1\par}
  \end{tcolorbox}%
  \nointerlineskip\franjas[4pt]
  \vspace{0.8em}
  \begin{tcolorbox}[enhanced, sharp corners, boxrule=0pt, frame hidden,
      colback=tonoPalido, borderline west={3pt}{0pt}{tonoPrincipal},
      borderline south={0.6pt}{0pt}{tonoClaro},
      left=10pt, right=8pt, top=6pt, bottom=6pt, before skip=0pt, after skip=0pt,
      fontupper=\small\color{tonoGrafito}]
    #2
  \end{tcolorbox}%
  \par\medskip}

%================================================================== documento ==
\begin{document}

\begin{titlepage}
\begin{tikzpicture}[remember picture,overlay]
  % banda superior
  \fill[tonoOscuro] (current page.north west) rectangle ([yshift=-5.4cm]current page.north east);
  \fill[tonoClaro] (current page.north west) rectangle ([xshift=1.1cm,yshift=-5.4cm]current page.north west);
  % franja de la gama bajo la banda
  \foreach \c [count=\i from 0] in {tonoOscuro,tonoProfundo,tonoPrincipal,tonoA,tonoB,tonoC,tonoClaro,tonoHielo}
    \fill[\c] ([xshift=\i*\paperwidth/8,yshift=-5.4cm]current page.north west)
      rectangle ++(\paperwidth/8,-0.35cm);
  % franja inferior
  \foreach \c [count=\i from 0] in {tonoHielo,tonoClaro,tonoC,tonoB,tonoA,tonoPrincipal,tonoProfundo,tonoOscuro}
    \fill[\c] ([xshift=\i*\paperwidth/8]current page.south west)
      rectangle ++(\paperwidth/8,0.6cm);
\end{tikzpicture}%
\centering
\vspace*{0.5cm}
{\sffamily\Huge\bfseries\color{white} Sílabo de Microeconometría\par}
\vspace{0.4cm}
{\sffamily\Large\color{tonoHielo} Plan de estudio integrado\par}
\vspace{1.9cm}
{\sffamily\large\bfseries\color{tonoPrincipal}\textcolor{tonoClaro}{\rule[0.5ex]{1.5cm}{0.6pt}}\quad Textos base\quad\textcolor{tonoClaro}{\rule[0.5ex]{1.5cm}{0.6pt}}\par}
\vspace{0.3cm}
{\large\etiqueta{colorCT}{CT} \quad {\color{tonoOscuro}A. Colin Cameron y Pravin K. Trivedi}\par}
{\normalsize \emph{\color{colorCT}Microeconometrics: Methods and Applications}\par}
{\small\color{tonoPizarra} Cambridge University Press, 2005 --- 27 capítulos\par}
\vspace{0.15cm}
{\large\etiqueta{colorCI}{CI} \quad {\color{tonoOscuro}Scott Cunningham}\par}
{\normalsize \emph{\color{colorCI}Causal Inference: The Remix}\par}
{\small\color{tonoPizarra} 2.ª edición en línea, \url{https://mixtape.scunning.com} --- 11 capítulos\par}
\vspace{0.45cm}
{\sffamily\large\bfseries\color{tonoPrincipal}\textcolor{tonoClaro}{\rule[0.5ex]{1.5cm}{0.6pt}}\quad Textos de apoyo\quad\textcolor{tonoClaro}{\rule[0.5ex]{1.5cm}{0.6pt}}\par}
\vspace{0.3cm}
{\large\etiqueta{colorMHE}{MHE} \quad {\color{tonoOscuro}Joshua D. Angrist y Jörn-Steffen Pischke}\par}
{\normalsize \emph{\color{colorMHE}Mostly Harmless Econometrics: An Empiricist's Companion}\par}
{\small\color{tonoPizarra} Princeton University Press, 2009 --- 8 capítulos\par}
\vspace{0.12cm}
{\large\etiqueta{colorMM}{MM} \quad {\color{tonoOscuro}Joshua D. Angrist y Jörn-Steffen Pischke}\par}
{\normalsize \emph{\color{colorMM}Mastering 'Metrics: The Path from Cause to Effect}\par}
{\small\color{tonoPizarra} Princeton University Press, 2015 --- 6 capítulos\par}
\vspace{0.12cm}
{\large\etiqueta{colorWd}{Wd} \quad {\color{tonoOscuro}Jeffrey M. Wooldridge}\par}
{\normalsize \emph{\color{colorWd}Econometric Analysis of Cross Section and Panel Data}\par}
{\small\color{tonoPizarra} MIT Press, 2002 --- 20 capítulos\par}
\vspace{0.12cm}
{\large\etiqueta{colorBt}{Bt} \quad {\color{tonoOscuro}Badi H. Baltagi}\par}
{\normalsize \emph{\color{colorBt}Econometric Analysis of Panel Data}\par}
{\small\color{tonoPizarra} Springer, 6.ª edición, 2021 --- 13 capítulos\par}
\vspace{0.12cm}
{\large\etiqueta{colorAr}{Ar} \quad {\color{tonoOscuro}Manuel Arellano}\par}
{\normalsize \emph{\color{colorAr}Panel Data Econometrics}\par}
{\small\color{tonoPizarra} Oxford University Press, 2003 --- 8 capítulos y 2 apéndices\par}
\vspace{0.12cm}
{\large\etiqueta{colorHu}{Hu} \quad {\color{tonoOscuro}Martin Huber}\par}
{\normalsize \emph{\color{colorHu}Causal Analysis: Impact Evaluation and Causal Machine Learning}\par}
{\small\color{tonoPizarra} MIT Press, 2023 --- 12 capítulos\par}
\vspace{0.35cm}
{\normalsize\color{colorPUCP} Con temas y lecturas de seis cursos de la PUCP\par}
\vfill
{\setlength{\fboxsep}{6pt}\colorbox{tonoPrincipal}{\normalsize\sffamily\bfseries\color{white} Trece bloques temáticos y un pretramo --- 7 meses y medio}\par}
\vspace{0.2cm}
\end{titlepage}

\tableofcontents
\newpage

%------------------------------------------------------------------------------
\section*{Cómo usar este sílabo}
\addcontentsline{toc}{section}{Cómo usar este sílabo}

El plan está organizado en \textbf{trece bloques temáticos} y un pretramo, y
sigue la estructura del curso 1ECO25 de J.~M.~Rentería: primero la inferencia
causal y los diseños de investigación (experimentos, selección en observables,
variables instrumentales y regresión discontinua), luego los datos de panel y
los diseños causales de panel (diferencias en diferencias, estudios de eventos
y control sintético), después los modelos de variable dependiente limitada y,
al final, los problemas de los datos: muestras complejas, error de medida y
datos faltantes. Conviene cerrar cada bloque antes de abrir el siguiente.

\medskip
\noindent\destacado{Requisito.} El plan supone el sílabo de econometría:
mínimos cuadrados y asintótica (Bloques 2 a 4), variables instrumentales
(Bloque 8), GMM (Bloque 9) y máxima verosimilitud (Bloque 10). Los modelos de
variable dependiente limitada y las ecuaciones simultáneas que allí se ven
(Bloques 11 y 12) se dan por sabidos: aquí solo se estudian sus extensiones y
sus versiones con datos de panel.

\medskip
\noindent\destacado{Criterio general de reparto} (en el orden en que se lee cada tema):
\begin{itemize}[leftmargin=1.4em,itemsep=2pt]
  \item \textcolor{colorMM}{\textbf{\emph{Mastering 'Metrics}}} y
        \textcolor{colorCI}{\textbf{Cunningham}} aportan la intuición y los
        ejemplos. Se leen \emph{primero} en los Bloques 1 a 9. \emph{The
        Remix} es la lectura principal de la inferencia causal y de los
        diseños de panel modernos (Bloques 8 y 9).
  \item \textcolor{colorMHE}{\textbf{\emph{Mostly Harmless Econometrics}}}
        da la versión formal de los mismos métodos.
  \item \textcolor{colorCT}{\textbf{Cameron y Trivedi}} es el texto base del
        plan: el marco general de la microeconometría, los modelos no lineales
        y los capítulos de datos de panel, muestras complejas, error de medida
        y datos faltantes.
  \item \textcolor{colorWd}{\textbf{Wooldridge}} es la alternativa rigurosa
        en datos de panel y en variable dependiente limitada.
  \item \textcolor{colorBt}{\textbf{Baltagi}} y
        \textcolor{colorAr}{\textbf{Arellano}} sostienen los Bloques 7 y 12
        y son la consulta de referencia del Bloque 6: Baltagi para el
        panorama completo y Arellano para la heterogeneidad no observada y
        los paneles dinámicos.
  \item \textcolor{colorHu}{\textbf{Huber}} cubre la evaluación de impacto y
        el aprendizaje automático causal, con aplicaciones en R.
\end{itemize}

\medskip
\noindent\destacado{Qué entra en las horas.} Las líneas \emph{Lecturas},
\emph{Profundización}, \emph{Artículos} y \emph{Artículos (PUCP)} entran en el
cómputo de horas. Las líneas \emph{Artículos} remiten a trabajos que no
figuran en los programas de la PUCP.
Las líneas \emph{Consulta} y \emph{Consulta (PUCP)} remiten a libros de la
carpeta y a referencias de los programas que no entran (Anexo E).

\medskip
\noindent\destacado{Material excluido.} Los capítulos de métodos generales de
Cameron y Trivedi (máxima verosimilitud, GMM, pruebas, métodos
semiparamétricos, bootstrap teórico y métodos bayesianos) ya están en el
sílabo de econometría, y lo mismo ocurre con los modelos básicos de variable
dependiente limitada (logit, probit, multinomial, ordenado, conteo, tobit y
Heckman, Bloque 12 de ese sílabo): los Bloques 10 y 11 solo tratan sus
extensiones. Los paneles no estacionarios y espaciales, los sistemas de
ecuaciones con componentes de error, los errores PCSE y la estimación
estructural pasan al material opcional. El detalle está en el Anexo B.

\medskip
\noindent\destacado{Sobre \emph{The Remix}.} Es la segunda edición, en línea,
de \emph{Causal Inference: The Mixtape} (2021), el libro que cita 1ECO25. Su
numeración cambia: los capítulos que Rentería asigna de la primera edición
(3, 4 y 7) son, en la segunda, los capítulos 3, 4 y 7 (DAG, resultados
potenciales y variables instrumentales), y la parte de diseños de panel se
amplía a cuatro capítulos (8 a 11). Trae códigos en R, Stata y Python.

\medskip
\noindent\destacado{Sobre el software y la replicación.} La microeconometría
aplicada se aprende replicando. Cameron y Trivedi tienen la versión en Stata
que recomiendan 1ECO25 y 1ECO77, \emph{Microeconometrics Using Stata} (2.ª ed.,
2022); Huber, Adams (2021) y Croissant y Millo (2019) trabajan en R. Los
paquetes de cada método están en el Anexo D.6. Conviene replicar cada semana un resultado publicado, como pide
el curso ECO315.

\medskip
\noindent\destacado{Notación.} \CT{n} remite al capítulo $n$ de Cameron y
Trivedi, y \CI{n}, \MHE{n}, \MM{n}, \WD{n}, \BT{n}, \AR{n} y \HU{n}, al
capítulo de los demás textos. Si solo entra parte de un capítulo se indican
secciones: \CT{21.6} es la sección 6 del capítulo 21; \MHE{3.3.2}, la
subsección 3.3.2. \emph{The Remix} se cita por capítulo.

%------------------------------------------------------------------------------
\newpage
\section*{Cronograma por semanas}
\addcontentsline{toc}{section}{Cronograma por semanas}

Calendario del plan: 320 horas distribuidas en \textbf{32 semanas de
contenido}, que equivalen a \textbf{unos 7 meses y medio de calendario}, a un
ritmo de diez horas semanales.

\medskip
\noindent Todas las semanas son de materia nueva y no hay semanas de colchón.
La columna de horas es la que gobierna el reparto. Si hay semanas perdidas por
feriados, viajes o imprevistos, el ajuste correcto es alargar el calendario,
no comprimir los bloques.

\medskip
{\filasalternas
\noindent\begin{tabularx}{\textwidth}{lXcc}
\toprule
\rowcolor{tonoOscuro}\cab{Semanas} & \cab{Contenido} & \cab{Sem.} & \cab{Horas} \\
\midrule
1 & Pretramo --- Datos microeconómicos y repaso & 1 & 10 \\
2--3 & Bloque 1 --- Causalidad: resultados potenciales y grafos & 2 & 20 \\
4--5 & Bloque 2 --- Experimentos aleatorios & 2 & 20 \\
\rowcolor{tonoHielo}\textbf{6--9} & \textbf{Bloque 3 --- Selección en observables} & \textbf{4} & \textbf{40} \\
\rowcolor{tonoHielo}\textbf{10--12} & \textbf{Bloque 4 --- Variables instrumentales y LATE} & \textbf{3} & \textbf{30} \\
13--14 & Bloque 5 --- Regresión discontinua & 2 & 20 \\
15--16 & Bloque 6 --- Datos de panel lineales & 2 & 20 \\
17--18 & Bloque 7 --- Paneles dinámicos & 2 & 20 \\
\rowcolor{tonoHielo}\textbf{19--22} & \textbf{Bloque 8 --- Diferencias en diferencias y estudios de eventos} & \textbf{4} & \textbf{40} \\
23--24 & Bloque 9 --- Control sintético y cortes transversales repetidos & 2 & 20 \\
25--26 & Bloque 10 --- Modelos de elección discreta & 2 & 20 \\
27--28 & Bloque 11 --- Censura, selección, conteo y duración & 2 & 20 \\
29--30 & Bloque 12 --- Modelos no lineales con datos de panel & 2 & 20 \\
31--32 & Bloque 13 --- Muestras complejas, error de medida y datos faltantes & 2 & 20 \\
\midrule
\rowcolor{tonoHielo}\multicolumn{2}{l}{\textbf{Total}} & \textbf{32} & \textbf{320} \\
\bottomrule
\end{tabularx}
}

\medskip
\noindent\destacado{Ritmo supuesto.} Diez horas por semana, todas de contenido
nuevo: por ejemplo, cuatro sesiones de dos horas entre semana para la teoría y
una sesión de dos horas el sábado para replicar un resultado con datos.

\medskip
\noindent\destacado{Los tramos largos.} Los Bloques 1 a 5 (semanas 2--14) forman
el recorrido por los diseños de identificación de corte transversal. Los
Bloques 6 a 9 (semanas 15--24) forman un recorrido continuo por los datos de
panel, de los estimadores clásicos a los diseños causales recientes, y son el
punto donde el avance suele atascarse.

\medskip
\noindent\destacado{Hitos de control.} Conviene verificar el avance al terminar las
semanas 9, 16, 24 y 28. Si en alguno de esos cortes se acumula más de una
semana de retraso, se alarga el calendario y no se comprime el bloque
siguiente.

\medskip
\noindent\destacado{Cambios respecto de la versión anterior.} Las semanas que se
liberan al no repetir el Bloque 12 del sílabo de econometría (Bloques 10 y 11
pasan de 3 a 2 semanas) y al recortar el panel clásico (Bloque 6 pasa de 3 a
2) se reasignan a la selección en observables (Bloque 3, de 3 a 4), a las
diferencias en diferencias (Bloque 8, de 3 a 4) y al control sintético
(Bloque 9, de 1 a 2). El total se mantiene en 32 semanas. Entran además las
pruebas múltiples y la inferencia por aleatorización (2.2), la sensibilidad
de Oster y de Cinelli--Hazlett (3.3), las descomposiciones (3.4), los
instrumentos débiles con las herramientas actuales (4.3), los instrumentos
\emph{shift-share} (4.4), la guía de Roodman para paneles dinámicos (7.3), los
estimadores robustos de diferencias en diferencias y las tendencias paralelas
aproximadas (8.4--8.6), el DiD sintético (9.2) y una tabla de software por
método (Anexo D.6).

\newpage
\subsection*{Detalle semana a semana}

{\small\filasalternas
\renewcommand{\arraystretch}{1.1}\noindent\begin{tabularx}{\textwidth}{clXc}
\toprule
\rowcolor{tonoOscuro}\cab{Sem.} & \cab{Bloque} & \cab{Temas} & \cab{Horas} \\
\midrule
1  & Pretramo & P.1 Datos microeconómicos; P.2 Repaso: MCO, variables instrumentales y máxima verosimilitud & 10 \\
\midrule
2  & 1 & 1.1 Modelos causales e identificación; 1.2 Resultados potenciales & 10 \\
3  & 1 & 1.3 Grafos acíclicos dirigidos y malos controles; 1.4 Regresión y causalidad & 10 \\
\midrule
4  & 2 & 2.1 Asignación aleatoria y análisis de experimentos; 2.2 Inferencia por aleatorización y pruebas múltiples & 10 \\
5  & 2 & 2.3 Experimentos en la práctica; 2.4 Experimentos naturales & 10 \\
\midrule
6  & 3 & 3.1 Independencia condicional, regresión y emparejamiento & 10 \\
7  & 3 & 3.2 Puntaje de propensión: emparejamiento y ponderación & 10 \\
8  & 3 & 3.3 Sensibilidad y cotas; 3.4 Descomposiciones: Oaxaca--Blinder, reponderación y RIF & 10 \\
9  & 3 & 3.5 Aprendizaje automático causal & 10 \\
\midrule
10 & 4 & 4.1 Estimador de Wald y MC2E & 10 \\
11 & 4 & 4.2 LATE y efectos heterogéneos & 10 \\
12 & 4 & 4.3 Instrumentos débiles e inferencia robusta; 4.4 Instrumentos \emph{shift-share} & 10 \\
\midrule
13 & 5 & 5.1 Regresión discontinua aguda; 5.2 Regresión discontinua difusa & 10 \\
14 & 5 & 5.3 Práctica de la regresión discontinua; 5.4 Diseños de quiebre y de aglomeración & 10 \\
\midrule
15 & 6 & 6.1 Datos de panel y heterogeneidad no observada; 6.2 Efectos fijos & 10 \\
16 & 6 & 6.3 Efectos aleatorios, Mundlak y Hausman; 6.4 Efectos de dos vías, inferencia y paneles no balanceados & 10 \\
\midrule
17 & 7 & 7.1 Variables instrumentales y GMM en panel; 7.2 Sesgo de Nickell & 10 \\
18 & 7 & 7.3 Arellano--Bond, Arellano--Bover y Blundell--Bond; 7.4 Variables predeterminadas & 10 \\
\midrule
19 & 8 & 8.1 Diseños causales de panel y DiD canónico; 8.2 DiD en regresión e inferencia & 10 \\
20 & 8 & 8.3 Estudios de eventos y tendencias previas; 8.4 Tendencias paralelas aproximadas & 10 \\
21 & 8 & 8.5 Adopción escalonada: por qué falla el estimador de dos vías & 10 \\
22 & 8 & 8.6 Estimadores robustos a efectos heterogéneos; 8.7 DiD doblemente robusto y triples diferencias & 10 \\
\midrule
23 & 9 & 9.1 Control sintético: estimación, requisitos e inferencia & 10 \\
24 & 9 & 9.2 Diferencias en diferencias sintéticas; 9.3 Pseudo paneles y cortes transversales repetidos & 10 \\
\midrule
25 & 10 & 10.1 Logit condicional y utilidad aleatoria; 10.2 Logit anidado, probit multinomial y logit mixto & 10 \\
26 & 10 & 10.3 Estimación por simulación; 10.4 Funciones de control en modelos no lineales & 10 \\
\midrule
27 & 11 & 11.1 Soluciones de esquina y modelo de dos partes; 11.2 Selección muestral y modelo de Roy & 10 \\
28 & 11 & 11.3 Extensiones de los modelos de conteo; 11.4 Modelos de duración & 10 \\
\midrule
29 & 12 & 12.1 Respuesta binaria con datos de panel & 10 \\
30 & 12 & 12.2 Tobit, selección y atrición en panel; 12.3 Conteo y modelos dinámicos no lineales & 10 \\
\midrule
31 & 13 & 13.1 Muestras estratificadas, ponderación y conglomerados; 13.2 Error de medida & 10 \\
32 & 13 & 13.3 Datos faltantes e imputación; 13.4 Bootstrap e inferencia en microeconometría & 10 \\
\bottomrule
\end{tabularx}
}

%------------------------------------------------------------------------------
\newpage
\bloque{Pretramo --- Datos microeconómicos y repaso --- 1 semana}
{Presenta las estructuras de datos de la microeconometría y repasa, en el
lenguaje de Cameron y Trivedi, los estimadores del sílabo de econometría que
se usan desde el Bloque 1.}

\tema{P.1 Datos microeconómicos: observacionales, experimentales y naturales}
\lecturas{\CT{1} \quad \CT{3.1--3.2} \quad \CT{3.5}}
\consultalib{Gertler et al. (2017), cap.~16}

\tema{P.2 Repaso: MCO, variables instrumentales y máxima verosimilitud}
\lecturas{\CT{4.1--4.9}}
\consultalib{Cameron y Trivedi (2005), caps.~5--6}
\nota{Equivale a los Bloques 2 a 4, 8 y 10 del sílabo de econometría. Si ese
plan ya se completó, la semana puede dedicarse a instalar el software y
preparar las bases de datos de los libros.}

%------------------------------------------------------------------------------
\bloque{Bloque 1 --- Causalidad: resultados potenciales y grafos --- 2 semanas}
{Los temas 4 de 1ECO25 (Rentería) y 4.1--4.2 de 1ECO25 (García), y la sesión
1.1 de 1ECO77: qué significa estimar un efecto causal y qué supuestos lo
permiten. Los dos lenguajes ---resultados potenciales y grafos--- se usan en el
resto del plan.}

\tema{1.1 Modelos causales, modelos estructurales e identificación}
\lecturas{\CI{1--2} \quad \CT{2.1--2.6} \quad \CT{2.8}}
\apoyo{\HU{1--2}}
\consulta{Imbens (2022) \quad Heckman y Pinto (2024)}

\tema{1.2 Resultados potenciales y el problema de selección}
\lecturas{\CI{4} \quad \MHE{2.1} \quad \CT{2.7} \quad \CT{25.1--25.3}}
\apoyo{\WD{18.1--18.2}}
\consultalib{Neal (2020), cap.~2}
\consulta{Rubin (1974) \quad Holland (1986)}
\nota{Efectos individuales, ATE, ATT y sesgo de selección.}

\tema{1.3 Grafos acíclicos dirigidos, buenos y malos controles}
\lecturas{\CI{3} \quad \MHE{3.2.3}}
\consultalib{Neal (2020), caps.~3--4 y 6}
\articulos{Cinelli, Forney y Pearl (2022)}
\consulta{Imbens (2020)}
\nota{Caminos por la puerta trasera, colisionadores y el criterio de ajuste.}

\tema{1.4 Regresión y causalidad}
\lecturas{\MM{2} \quad \MHE{3.1--3.2}}
\nota{La regresión como aproximación a la esperanza condicional, la fórmula
del sesgo de variable omitida y los modelos saturados.}

%------------------------------------------------------------------------------
\bloque{Bloque 2 --- Experimentos aleatorios --- 2 semanas}
{El tema 5 de 1ECO25 (Rentería) y el tema 4.3 de 1ECO25 (García). El
experimento es la referencia contra la que se juzgan los demás diseños.}

\tema{2.1 Asignación aleatoria y análisis de experimentos}
\lecturas{\MM{1} \quad \MHE{2.2--2.3} \quad \CT{3.3}}
\apoyo{\HU{3}}
\articulosB{Athey e Imbens (2017)}
\consultalib{Gertler et al. (2017), cap.~4 \quad Bailey (2020), cap.~10}
\nota{Athey e Imbens presentan el análisis de experimentos desde la propia
aleatorización: el estimador de Neyman, la regresión con covariables y los
diseños estratificados y por conglomerados.}

\tema{2.2 Inferencia por aleatorización y pruebas múltiples}
\articulosB{Young (2019) \quad Anderson (2008) \quad List, Shaikh y Xu (2019)}
\nota{Pruebas exactas de Fisher y valores $p$ por permutación, y por qué
cambian las conclusiones de muchos experimentos publicados (Young, 2019).
Corrección por pruebas múltiples: control de la tasa de error por familia
(Bonferroni, Holm, Romano--Wolf) y de la tasa de falsos descubrimientos, e
índices resumen de resultados (Anderson, 2008).}

\tema{2.3 Experimentos en la práctica: diseño, potencia, estratificación, cumplimiento y atrición}
\lecturas{Gertler et al. (2017), caps.~9 y 15}
\articulos{Duflo, Glennerster y Kremer (2006)}
\articulosB{Bruhn y McKenzie (2009)}
\consulta{Glennerster y Takavarasha (2013) \quad Bédécarrats, Guérin y Roubaud (2020)}
\nota{Tamaño de muestra y cálculo de potencia, aleatorización por
conglomerados, estratificación y emparejamiento por pares para mejorar el
balance (Bruhn y McKenzie, 2009), cumplimiento imperfecto (que se retoma en el
tema 4.2) y pérdida de muestra.}

\tema{2.4 Experimentos naturales}
\lecturas{\CT{3.4}}
\articulos{Dammert, Galdo y Galdo (2014)}
\nota{Un experimento de campo en el Perú sobre prevención del dengue con
mensajes de texto (lectura de 1ECO77).}

%------------------------------------------------------------------------------
\bloque{Bloque 3 --- Selección en observables --- 4 semanas}
{El tema 6 de 1ECO25 (Rentería), el tema 4.4 de 1ECO25 (García) y el módulo 1
de 1ECO77: cuándo controlar por características observables basta para
identificar un efecto causal y cómo hacerlo. Incluye las descomposiciones de
brechas, que descansan en el mismo supuesto, y dedica una semana completa al
aprendizaje automático causal.}

\tema{3.1 Independencia condicional, regresión y emparejamiento}
\lecturas{\CI{5} \quad \MHE{3.2.1} \quad \MHE{3.3.1}}
\apoyo{\WD{18.3}}
\consultalib{García Núñez (2010)}

\tema{3.2 El puntaje de propensión: emparejamiento, estratificación y ponderación}
\lecturas{\MHE{3.3.2--3.3.3} \quad \CT{25.4} \quad \HU{4}}
\consultalib{Gertler et al. (2017), cap.~8}
\articulos{Rosenbaum y Rubin (1983) \quad Dehejia y Wahba (2002)}
\consulta{Caliendo y Kopeinig (2008) \quad Abadie e Imbens (2016) \quad LaLonde (1986) \quad Smith y Todd (2005)}
\nota{Soporte común, balance de covariables, ponderación por la inversa de la
probabilidad (IPW) y estimadores doblemente robustos.}

\tema{3.3 Sensibilidad y cotas}
\lecturas{\HU{10}}
\articulosB{Oster (2019) \quad Cinelli y Hazlett (2020)}
\consultalib{Neal (2020), cap.~8}
\nota{Qué se puede concluir si el supuesto de independencia condicional
falla: cotas de Manski y análisis de sensibilidad. Oster (2019) acota el
sesgo por no observables a partir de cómo cambian el coeficiente y el $R^2$ al
añadir controles; Cinelli y Hazlett (2020) lo expresan con la fórmula del sesgo
de variable omitida y el valor de robustez (paquete \texttt{sensemakr}). Son
los dos análisis de sensibilidad que más se piden hoy en los artículos
aplicados.}

\tema{3.4 Descomposiciones: Oaxaca--Blinder, reponderación y regresiones RIF}
\articulosB{Fortin, Lemieux y Firpo (2011), secc.~1--3 \quad Firpo, Fortin y Lemieux (2009)}
\consultalib{Oaxaca (1973) \quad Blinder (1973) \quad DiNardo, Fortin y Lemieux (1996) \quad Fortin, Lemieux y Firpo (2011), secc.~4--5}
\nota{La descomposición de Oaxaca--Blinder de una brecha de medias en una parte
explicada por las características y otra por los retornos, y su lectura como
un efecto de tratamiento bajo independencia condicional; la reponderación de
DiNardo, Fortin y Lemieux, que usa el puntaje de propensión del tema 3.2, y
las regresiones de la función de influencia recentrada (RIF) para descomponer
cuantiles, varianzas o índices de Gini. Aplicación natural con la ENAHO:
brechas salariales por género, entre empleo formal e informal o entre Lima y
el resto del país.}

\tema{3.5 Aprendizaje automático causal}
\lecturas{\HU{5}}
\consultalib{Hansen (2021), caps.~29.20--29.22}
\articulos{Chernozhukov et al. (2018)}
\nota{Aprendizaje automático doble/insesgado y efectos heterogéneos. Retoma el
tema 13.6 del sílabo de econometría y el curso ECO386 (Anexo D.3). Ocupa una
semana completa porque el artículo de Chernozhukov et al. es denso y conviene
implementarlo.}

%------------------------------------------------------------------------------
\bloque{Bloque 4 --- Variables instrumentales y LATE --- 3 semanas}
{El tema 10 de 1ECO25 (Rentería), el tema 4.7 de 1ECO25 (García) y el módulo 2
de 1ECO77. Supone el Bloque 8 del sílabo de econometría; aquí el énfasis está
en la interpretación causal con efectos heterogéneos.}

\tema{4.1 El estimador de Wald y MC2E}
\lecturas{\MM{3} \quad \CI{7} \quad \MHE{4.1--4.3}}
\consultalib{\MM{6} \quad Bailey (2020), cap.~9}

\tema{4.2 Efecto promedio local del tratamiento y efectos heterogéneos}
\lecturas{\MHE{4.4--4.5} \quad \CT{25.7} \quad \HU{6}}
\apoyo{\WD{18.4}}
\articulos{Angrist (2022)}
\consulta{Mogstad y Torgovitsky (2024)}
\nota{Cumplidores, siempre tomadores y nunca tomadores; variables
instrumentales en experimentos con cumplimiento imperfecto.}

\tema{4.3 Instrumentos débiles e inferencia robusta}
\lecturas{\MHE{4.6.4} \quad \CT{4.9}}
\articulosB{Andrews, Stock y Sun (2019) \quad Lee, McCrary, Moreira y Porter (2022)}
\consultalib{Hansen (2021), caps.~12.35--12.39}
\consulta{Anderson y Rubin (1949) \quad Stock y Yogo (2005) \quad Montiel Olea y Pflueger (2013)}
\nota{El estadístico $F$ efectivo de Montiel Olea y Pflueger como diagnóstico
con errores robustos o por conglomerados; la inferencia de Anderson--Rubin,
robusta a instrumentos débiles, y la corrección del estadístico $t$ de Lee et
al. (el procedimiento $tF$). Las pruebas de Cragg--Donald y de Stock--Yogo
suponen errores homoscedásticos y quedan como referencia. Aplicación a
Angrist y Krueger (1991).}

\tema{4.4 Instrumentos de participación y desplazamiento (\emph{shift-share} o Bartik)}
\articulosB{Goldsmith-Pinkham, Sorkin y Swift (2020) \quad Borusyak, Hull y Jaravel (2022)}
\consultalib{Cai (2024), sección de instrumentos con exposición a choques}
\nota{Instrumentos construidos como la exposición de cada unidad
(participaciones iniciales) a choques comunes (desplazamientos). Hay dos
formas de justificar la identificación: por la exogeneidad de las
participaciones (Goldsmith-Pinkham et al.) o por la de los choques (Borusyak
et al.), con inferencias distintas en cada caso. Aplicación natural en el
Perú: la exposición regional a los precios de los minerales o a choques de
comercio exterior. Las funciones de control y la endogeneidad en modelos no
lineales se estudian en el tema 10.4, una vez vistos los modelos de elección
discreta.}

%------------------------------------------------------------------------------
\bloque{Bloque 5 --- Regresión discontinua --- 2 semanas}
{El tema 11 de 1ECO25 (Rentería) y el tema 4.6 de 1ECO25 (García): diseños
aguda y difusa con las herramientas de inferencia actuales.}

\tema{5.1 Regresión discontinua aguda}
\lecturas{\MM{4} \quad \CI{6} \quad \MHE{6.1}}
\apoyo{\CT{25.6}}

\tema{5.2 Regresión discontinua difusa}
\lecturas{\MHE{6.2}}
\apoyo{\HU{9}}
\nota{La regresión discontinua difusa como un caso de variables
instrumentales.}

\tema{5.3 Práctica de la regresión discontinua: ancho de banda, inferencia y validación}
\articulos{Cattaneo, Idrobo y Titiunik (2020) \quad Cattaneo y Titiunik (2022)}
\consultalib{Hansen (2021), cap.~21}
\consulta{Imbens y Lemieux (2007)}
\consultalib{Cattaneo, Idrobo y Titiunik (2024) \quad Gelman e Imbens (2019)}
\nota{Polinomios locales, ancho de banda óptimo, inferencia robusta con
corrección de sesgo, pruebas de manipulación y de balance.}
\nota{El volumen de 2020 de Cattaneo et al. trabaja el marco basado en
continuidad; el de 2024 añade el marco de aleatorización local, los diseños
difusos y las variables de asignación discretas. Gelman e Imbens (2019) explican
por qué no conviene usar polinomios globales de orden alto. Los dos marcos y los
polinomios de orden alto son tema del curso de Econometría del BCRP.}

\tema{5.4 Diseños de quiebre y de aglomeración}
\lecturas{\HU{9}}
\consultalib{Cai (2024), secciones de RD y \emph{bunching}}
\nota{Los diseños de quiebre (\emph{kink}) y de aglomeración
(\emph{bunching}) como extensiones de la regresión discontinua.}

%------------------------------------------------------------------------------
\bloque{Bloque 6 --- Datos de panel lineales --- 2 semanas}
{El tema 7 y el tema 12 de 1ECO25 (Rentería), el tema 3 de 1ECO25 (García) y
el módulo 3 de 1ECO77. Es el bloque nuevo respecto del sílabo de
econometría, que dejó fuera los datos de panel. Se apoya en los libros de la
carpeta Panel Data. Cada tema tiene una sola lectura principal; los demás
textos quedan como consulta. Los sistemas de ecuaciones con componentes de
error y los errores PCSE pasan al Anexo D.1, y las variables instrumentales en
panel abren el Bloque 7.}

\tema{6.1 Datos de panel y heterogeneidad no observada}
\lecturas{\CT{21.1--21.5}}
\apoyo{\AR{1} \quad \AR{2.1}}
\consultalib{\WD{10.1--10.3} \quad \BT{1} \quad Hsiao (2003), cap.~1 \quad Biørn (2017), cap.~1}
\nota{Ventajas y límites de los datos de panel; el modelo agrupado.}

\tema{6.2 Efectos fijos: estimador intragrupos, variables ficticias y primeras diferencias}
\lecturas{\CT{21.6} \quad \CI{8}}
\apoyo{\AR{2.2}}
\consultalib{\WD{10.5--10.6} \quad \MHE{5.1} \quad \BT{2.1--2.2} \quad Bailey (2020), cap.~8 \quad Wooldridge (2012), caps.~13--14}
\consulta{Imai y Kim (2019)}
\nota{Estimadores \emph{within}, LSDV, \emph{between} y en primeras
diferencias; qué identifica cada uno. \CI{8} da la lectura causal de los
efectos fijos que retoma el Bloque 8.}

\tema{6.3 Efectos aleatorios, efectos aleatorios correlacionados y la prueba de Hausman}
\lecturas{\CT{21.7}}
\apoyo{\WD{11.4}}
\consultalib{\WD{10.4} \quad \WD{10.7} \quad \BT{2.3--2.4} \quad \BT{4.2--4.3} \quad \BT{7.4} \quad \AR{3.1--3.4} \quad Beltrán y Castro (2010), cap.~2 \quad Hsiao (2003), cap.~3.9}
\nota{El estimador de MCG de Balestra--Nerlove, la prueba de Breusch--Pagan y
la prueba de Hausman; los enfoques de Mundlak y de Chamberlain para relajar el
supuesto de efectos aleatorios. El estimador de Hausman--Taylor para variables
constantes en el tiempo (\BT{7.4}) queda como consulta.}

\tema{6.4 Efectos de dos vías, inferencia y paneles no balanceados}
\lecturas{\CT{21.8}}
\apoyo{\BT{3}}
\consultalib{\BT{5} \quad \BT{9.1--9.2} \quad Croissant y Millo (2019)}
\nota{Efectos fijos de individuo y de tiempo; heteroscedasticidad, correlación
serial y errores agrupados por individuo; paneles no balanceados y por qué
faltan observaciones. \BT{4.1} trata las pruebas de agrupabilidad.}

%------------------------------------------------------------------------------
\bloque{Bloque 7 --- Paneles dinámicos --- 2 semanas}
{El tema 3.5 de 1ECO25 (García) y la sesión 4.1 de 1ECO77. Arellano es la
lectura principal. Abre con las variables instrumentales y el GMM en paneles
estáticos, la base de los estimadores de Arellano y Bond.}

\tema{7.1 Variables instrumentales y GMM con datos de panel}
\lecturas{\CT{22.1--22.4}}
\apoyo{\WD{11.1--11.3}}
\consultalib{\AR{4} \quad \BT{7.1--7.3}}
\nota{Endogeneidad en paneles estáticos: variables instrumentales con efectos
fijos, instrumentos internos y el estimador GMM de panel.}

\tema{7.2 El sesgo de Nickell}
\lecturas{\AR{6.1--6.3} \quad \BT{8.1}}
\apoyo{\CT{22.5}}
\consultalib{Hsiao (2003), cap.~4}

\tema{7.3 Los estimadores de Arellano--Bond, Arellano--Bover y Blundell--Bond}
\lecturas{\BT{8.2--8.5} \quad \AR{6.4--6.5}}
\articulosB{Roodman (2009a) \quad Roodman (2009b)}
\consultalib{Pesaran (2015), cap.~27 \quad Croissant y Millo (2019)}
\nota{GMM en diferencias y en sistema, pruebas de Sargan--Hansen y de
autocorrelación, y el problema de los instrumentos numerosos. Roodman (2009a)
es la guía práctica de \texttt{xtabond2}; Roodman (2009b) muestra cómo el
exceso de instrumentos sesga las estimaciones y debilita la prueba de Hansen,
y cómo limitarlo (colapsar la matriz de instrumentos y acotar los rezagos).}

\tema{7.4 Variables predeterminadas y efectos fijos frente a rezagos}
\lecturas{\AR{8} \quad \MHE{5.3}}
\consultalib{\AR{7} \quad \BT{8.8}}

%------------------------------------------------------------------------------
\bloque{Bloque 8 --- Diferencias en diferencias y estudios de eventos --- 4 semanas}
{Los temas 8 y 9 de 1ECO25 (Rentería), el tema 4.5 de 1ECO25 (García) y la
sesión 1.2 de 1ECO77. \emph{The Remix} es la lectura principal: sus capítulos
9 y 10 recogen la literatura reciente sobre efectos heterogéneos. Los
artículos que proponen los estimadores robustos entran en las horas, porque
son los que se usan en la práctica.}

\tema{8.1 Diseños causales de panel y diferencias en diferencias canónico}
\lecturas{\MM{5} \quad \CI{8--9} \quad \MHE{5.2}}
\apoyo{\CT{22.6} \quad \CT{25.5} \quad \HU{7}}
\consultalib{Gertler et al. (2017), cap.~7}
\articulos{Card y Krueger (1994) \quad Galdo (2013)}
\nota{Galdo (2013) estima con este diseño los efectos de largo plazo de la
exposición temprana a la violencia de Sendero Luminoso.}

\tema{8.2 Diferencias en diferencias en regresión e inferencia}
\lecturas{\MHE{5.2.1} \quad \MHE{8.2}}
\nota{Efectos fijos de dos vías, controles, correlación serial y pocos
conglomerados.}

\tema{8.3 Estudios de eventos y tendencias previas}
\lecturas{\CI{10}}
\articulos{Miller (2023) \quad Roth (2022)}
\consulta{Freyaldenhoven, Hansen, Pérez y Shapiro (2024)}
\nota{Especificación del estudio de eventos, normalización y ventanas; por qué
no basta con no rechazar la prueba de tendencias previas (Roth, 2022).}

\tema{8.4 Tendencias paralelas aproximadas e inferencia robusta}
\articulosB{Rambachan y Roth (2023)}
\nota{Cotas a la violación de las tendencias paralelas a partir de las
tendencias previas (restricciones de magnitud relativa y de suavidad) e
intervalos de confianza robustos con el paquete \texttt{HonestDiD}.}

\tema{8.5 Adopción escalonada: por qué falla el estimador de efectos fijos de dos vías}
\lecturas{\CI{10}}
\articulos{Goodman-Bacon (2021) \quad de Chaisemartin y D'Haultfœuille (2020) \quad Roth, Sant'Anna, Bilinski y Poe (2023)}
\nota{La descomposición de Goodman-Bacon, las comparaciones con unidades ya
tratadas y los pesos negativos; la síntesis de Roth et al. como mapa de la
literatura.}

\tema{8.6 Estimadores robustos a efectos heterogéneos}
\articulos{Callaway y Sant'Anna (2021) \quad Sun y Abraham (2021)}
\articulosB{Borusyak, Jaravel y Spiess (2024)}
\nota{Efectos promedio por grupo y período y su agregación
(Callaway--Sant'Anna), estudios de eventos con interacciones por cohorte
(Sun--Abraham) y el estimador por imputación (Borusyak et al.). Conviene
estimar los tres con los mismos datos y comparar los resultados.}

\tema{8.7 Diferencias en diferencias doblemente robusto, con emparejamiento y triples diferencias}
\lecturas{\CI{10}}
\articulos{Sant'Anna y Zhao (2020)}
\consultalib{García Núñez (2010)}

%------------------------------------------------------------------------------
\bloque{Bloque 9 --- Control sintético y cortes transversales repetidos --- 2 semanas}
{El control sintético del curso ECO315, su combinación con las diferencias en
diferencias y dos estructuras de datos frecuentes en el Perú, donde los
paneles largos son escasos.}

\tema{9.1 Control sintético: estimación, requisitos e inferencia}
\lecturas{\CI{11} \quad \HU{8}}
\articulosB{Abadie, Diamond y Hainmueller (2010) \quad Abadie (2021)}
\consultalib{Adams (2021), cap.~11}
\nota{Elección de los pesos y de los predictores, ajuste previo al
tratamiento, cuándo el método es factible y qué datos exige (Abadie, 2021), e
inferencia por placebos en el espacio y en el tiempo.}

\tema{9.2 Diferencias en diferencias sintéticas}
\articulosB{Arkhangelsky, Athey, Hirshberg, Imbens y Wager (2021)}
\nota{Pesos por unidad y por período que combinan el control sintético con
las diferencias en diferencias. Conviene comparar los tres estimadores con los
mismos datos (paquete \texttt{synthdid}).}

\tema{9.3 Pseudo paneles y cortes transversales repetidos}
\lecturas{\CT{22.7} \quad \BT{10.3}}
\consultalib{Wooldridge (2012), cap.~13}
\nota{Útil con encuestas como la ENAHO, que combina cortes transversales con un
componente de panel.}

%------------------------------------------------------------------------------
\bloque{Bloque 10 --- Modelos de elección discreta --- 2 semanas}
{El tema 3 de 1ECO25 (Rentería), el tema 2.1--2.2 de 1ECO25 (García) y la
Unidad 3 de 2ECO01. Supone el Bloque 12 del sílabo de econometría (modelo de
probabilidad lineal, logit, probit, logit multinomial, ordenado y anidado) y no
lo repite: aquí entran los modelos de utilidad aleatoria de Train, la
estimación por simulación y la endogeneidad en modelos no lineales. Los
capítulos de Cameron y Trivedi y de Wooldridge sobre los modelos básicos
quedan como consulta para repasar. El logit mixto y la simulación preparan la
estimación de demanda del sílabo de organización industrial (tema 12.2).}

\tema{10.1 Logit condicional y modelos de utilidad aleatoria}
\lecturas{\CT{15.1--15.5}}
\apoyo{Train (2009), caps.~2--3}
\consultalib{\CT{14} \quad \WD{15.1--15.6} \quad \WD{15.9} \quad Beltrán y Castro (2010), caps.~3--4 \quad Bailey (2020), cap.~12}
\nota{Del logit multinomial al condicional: atributos que varían por
alternativa, el supuesto de independencia de alternativas irrelevantes y cómo
probarlo. Las muestras basadas en la elección (\CT{14}) quedan como consulta.}

\tema{10.2 Logit anidado, probit multinomial y logit mixto}
\lecturas{\CT{15.6--15.8}}
\consultalib{Train (2009), caps.~4--6 \quad \CT{15.9--15.10} \quad \WD{15.10}}
\nota{Las respuestas ordenadas y multivariadas (\CT{15.9--15.10}) ya están en
el Bloque 12 del sílabo de econometría y quedan como consulta.}

\tema{10.3 Estimación por simulación}
\lecturas{Train (2009), caps.~9--10}
\consultalib{Train (2009), cap.~8 \quad Cameron y Trivedi (2005), cap.~12}
\nota{Extracción de densidades y máxima verosimilitud simulada. El cap.~8 de
Train (maximización numérica) repite el sílabo de econometría.}

\tema{10.4 Funciones de control y endogeneidad en modelos no lineales}
\lecturas{\MHE{4.6.3} \quad \WD{15.7} \quad \WD{19.5}}
\consulta{Chay y Greenstone (2005) \quad Evans y Schwab (1995)}
\nota{La función de control como alternativa a MC2E, el modelo de
coeficientes aleatorios correlacionados y el probit con regresor endógeno
(sesión 2.3 de 1ECO77). Retoma el Bloque 4.}

%------------------------------------------------------------------------------
\bloque{Bloque 11 --- Censura, selección, conteo y duración --- 2 semanas}
{El tema 13 de 1ECO25 (Rentería), los temas 2.3--2.4 de 1ECO25 (García) y los
modelos de duración de su sumilla. Supone el Bloque 12 del sílabo de
econometría (tobit, Heckman y modelos básicos de conteo) y se concentra en sus
extensiones; los capítulos de Wooldridge quedan como consulta.}

\tema{11.1 Soluciones de esquina, el modelo tobit y el modelo de dos partes}
\lecturas{\CT{16.1--16.4}}
\consultalib{\WD{16.1--16.7} \quad Beltrán y Castro (2010), cap.~5}
\nota{Repaso breve del tobit; el énfasis está en el modelo de dos partes y en
cuándo conviene frente al tobit.}

\tema{11.2 Selección muestral y el modelo de Roy}
\lecturas{\CT{16.5--16.9}}
\consultalib{\WD{17.1--17.6}}
\nota{El modelo de Heckman ya se vio en el sílabo de econometría; aquí entran
el modelo de Roy, la estimación semiparamétrica y su relación con el LATE del
tema 4.2.}

\tema{11.3 Extensiones de los modelos de conteo}
\lecturas{\CT{20}}
\consultalib{\WD{19.1--19.5}}
\nota{Sobredispersión y binomial negativa, modelos con exceso de ceros y de
valla (\emph{hurdle}), y el Poisson como estimador de cuasi-máxima
verosimilitud robusto.}

\tema{11.4 Modelos de duración}
\lecturas{\CT{17}}
\consultalib{Wooldridge (2002), cap.~20}

%------------------------------------------------------------------------------
\bloque{Bloque 12 --- Modelos no lineales con datos de panel --- 2 semanas}
{La sesión 4.2 de 1ECO77. Une los Bloques 6 y 7 con los Bloques 10 y 11.}

\tema{12.1 Respuesta binaria con datos de panel}
\lecturas{\CT{23.1--23.4} \quad \WD{15.8}}
\apoyo{\BT{11.1--11.2}}
\nota{Logit condicional de efectos fijos, probit de efectos aleatorios y
efectos aleatorios correlacionados.}

\tema{12.2 Tobit, selección y atrición con datos de panel}
\lecturas{\CT{23.5} \quad \WD{16.8} \quad \WD{17.7}}
\apoyo{\BT{11.4--11.5}}

\tema{12.3 Conteo y modelos dinámicos no lineales con datos de panel}
\lecturas{\CT{23.6--23.7} \quad \WD{19.6}}
\apoyo{\BT{10.6} \quad \BT{11.3}}

%------------------------------------------------------------------------------
\bloque{Bloque 13 --- Muestras complejas, error de medida y datos faltantes --- 2 semanas}
{Los problemas de los datos reales. Sigue el curso ECO387 (2022-2) en lo
relativo a encuestas como la ENAHO.}

\tema{13.1 Muestras estratificadas, ponderación y conglomerados}
\lecturas{\CT{24} \quad \WD{17.8}}
\articulos{Cameron y Miller (2015) \quad Solon, Haider y Wooldridge (2015)}
\consulta{Abadie, Athey, Imbens y Wooldridge (2017)}

\tema{13.2 Error de medida}
\lecturas{\CT{26} \quad \AR{4.1--4.2}}
\apoyo{\BT{10.1}}

\tema{13.3 Datos faltantes e imputación}
\lecturas{\CT{27}}

\tema{13.4 Bootstrap e inferencia en microeconometría}
\lecturas{\CT{11} \quad \MHE{8.1}}
\nota{Errores estándar robustos en muestras pequeñas y bootstrap en las
aplicaciones de este plan. La teoría está en el Bloque 5 del sílabo de
econometría.}

%------------------------------------------------------------------------------
\newpage
\section*{Anexo A. De los capítulos a los bloques}
\addcontentsline{toc}{section}{Anexo A. De los capítulos a los bloques}

Capítulos de cada libro que entran en el plan y temas donde se leen. Las
secciones excluidas se detallan en el Anexo B.

\newcommand{\mapa}[2]{%
\subsection*{#1}
{\small\filasalternas
\noindent\begin{tabularx}{\textwidth}{>{\raggedright\arraybackslash}p{1.1cm}X>{\raggedright\arraybackslash}p{3.8cm}}
\toprule
\rowcolor{tonoOscuro}\cab{Cap.} & \cab{Título} & \cab{Temas} \\
\midrule
#2
\bottomrule
\end{tabularx}
}}

\mapa{A.1 Cameron y Trivedi, \emph{Microeconometrics: Methods and Applications} (texto base)}{%
1 & Overview & P.1 \\
2 & Causal and Noncausal Models & 1.1, 1.2 \\
3 & Microeconomic Data Structures & P.1, 2.1, 2.4 \\
4 & Linear Models (4.1--4.9) & P.2, 4.3 \\
11 & Bootstrap Methods & 13.4 \\
14 & Binary Outcome Models (consulta) & 10.1 \\
15 & Multinomial Models (15.1--15.8; 15.9--15.10 consulta) & 10.1, 10.2 \\
16 & Tobit and Selection Models (16.1--16.9) & 11.1, 11.2 \\
17 & Transition Data: Survival Analysis & 11.4 \\
20 & Models of Count Data & 11.3 \\
21 & Linear Panel Models: Basics & 6.1, 6.2, 6.3, 6.4 \\
22 & Linear Panel Models: Extensions (22.1--22.7) & 7.1, 7.2, 8.1, 9.3 \\
23 & Nonlinear Panel Models (23.1--23.7) & 12.1, 12.2, 12.3 \\
24 & Stratified and Clustered Samples & 13.1 \\
25 & Treatment Evaluation & 1.2, 3.2, 4.2, 5.1, 8.1 \\
26 & Measurement Error Models & 13.2 \\
27 & Missing Data and Imputation & 13.3 \\
}

\mapa{A.2 Cunningham, \emph{Causal Inference: The Remix} (texto base)}{%
1 & Foundational Ideas in Causal Inference & 1.1 \\
2 & Which Causal Inference? & 1.1 \\
3 & Directed Acyclic Graphs & 1.3 \\
4 & Potential Outcomes and Randomization & 1.2 \\
5 & Unconfoundedness & 3.1 \\
6 & Regression Discontinuity & 5.1 \\
7 & Instrumental Variables & 4.1 \\
8 & Causal Panel Designs & 6.2, 8.1 \\
9 & Difference-in-Differences Fundamentals & 8.1 \\
10 & Complex Diff-in-Diff Designs & 8.3, 8.5, 8.7 \\
11 & Synthetic Control & 9.1 \\
}

\mapa{A.3 Angrist y Pischke, \emph{Mostly Harmless Econometrics}}{%
2 & The Experimental Ideal & 1.2, 2.1 \\
3 & Making Regression Make Sense (3.1--3.3) & 1.3, 1.4, 3.1, 3.2 \\
4 & Instrumental Variables in Action & 4.1, 4.2, 4.3, 10.4 \\
5 & Parallel Worlds: Fixed Effects, DiD, and Panel Data & 6.2 (consulta), 7.4, 8.1, 8.2 \\
6 & Getting a Little Jumpy: Regression Discontinuity & 5.1, 5.2 \\
8 & Nonstandard Standard Error Issues & 8.2, 13.4 \\
}

\mapa{A.4 Angrist y Pischke, \emph{Mastering 'Metrics}}{%
1 & Randomized Trials & 2.1 \\
2 & Regression & 1.4 \\
3 & Instrumental Variables & 4.1 \\
4 & Regression Discontinuity Designs & 5.1 \\
5 & Differences-in-Differences & 8.1 \\
}

\mapa{A.5 Wooldridge, \emph{Econometric Analysis of Cross Section and Panel Data}}{%
10 & Basic Linear Unobserved Effects Panel Data Models (consulta) & 6.1, 6.2, 6.3 \\
11 & More Topics in Linear Unobserved Effects Models (11.1--11.4) & 6.3, 7.1 \\
15 & Discrete Response Models (15.7--15.8; el resto, consulta) & 10.4, 12.1; consulta en 10.1, 10.2 \\
16 & Corner Solution Outcomes and Censored Regression (16.8; 16.1--16.7 consulta) & 12.2; consulta en 11.1 \\
17 & Sample Selection, Attrition, and Stratified Sampling (17.7--17.8; 17.1--17.6 consulta) & 12.2, 13.1; consulta en 11.2 \\
18 & Estimating Average Treatment Effects (18.1--18.4) & 1.2, 3.1, 4.2 \\
19 & Count Data and Related Models (19.5--19.6; 19.1--19.4 consulta) & 10.4, 12.3; consulta en 11.3 \\
}

\mapa{A.6 Baltagi, \emph{Econometric Analysis of Panel Data}}{%
1 & Introduction (consulta) & 6.1 \\
2 & The One-Way Error Component Regression Model (2.1--2.4, consulta) & 6.2, 6.3 \\
3 & The Two-Way Error Component Regression Model & 6.4 \\
4 & Test of Hypotheses with Panel Data (4.1--4.3, consulta) & 6.3, 6.4 \\
5 & Heteroskedasticity and Serial Correlation (consulta) & 6.4 \\
7 & Simultaneous Equations with Error Components (7.1--7.4, consulta) & 6.3, 7.1 \\
8 & Dynamic Panel Data Models (8.1--8.5; 8.8 consulta) & 7.2, 7.3, 7.4 \\
9 & Unbalanced Panel Data Models (9.1--9.2, consulta) & 6.4 \\
10 & Special Topics (10.1, 10.3, 10.6) & 9.3, 12.3, 13.2 \\
11 & Limited Dependent Variables and Panel Data (11.1--11.5) & 12.1, 12.2, 12.3 \\
}

\mapa{A.7 Arellano, \emph{Panel Data Econometrics}}{%
1 & Introduction & 6.1 \\
2 & Unobserved Heterogeneity (2.1--2.2) & 6.1, 6.2 \\
3 & Error Components (3.1--3.4, consulta) & 6.3 \\
4 & Error in Variables & 13.2; consulta en 7.1 \\
6 & Autoregressive Models with Individual Effects (6.1--6.5) & 7.2, 7.3 \\
7 & Models with both Strictly Exogenous and Lagged Dependent Variables (consulta) & 7.4 \\
8 & Predetermined Variables & 7.4 \\
}

\mapa{A.8 Huber, \emph{Causal Analysis}}{%
1--2 & Introduction; Causality and no causality & 1.1 \\
3 & Social experiments and linear regression & 2.1 \\
4 & Selection on observables & 3.2 \\
5 & Causal machine learning & 3.5 \\
6 & Instrumental variables & 4.2 \\
7 & Difference-in-Differences & 8.1 \\
8 & Synthetic controls & 9.1 \\
9 & Regression discontinuity, kink, and bunching designs & 5.2, 5.4 \\
10 & Partial identification and sensitivity analysis & 3.3 \\
}

%------------------------------------------------------------------------------
\section*{Anexo B. Capítulos y secciones que no entran en el plan}
\addcontentsline{toc}{section}{Anexo B. Capítulos y secciones que no entran en el plan}

{\small\filasalternas
\noindent\begin{tabularx}{\textwidth}{>{\raggedright\arraybackslash}p{4.6cm}X}
\toprule
\rowcolor{tonoOscuro}\cab{Materia} & \cab{Capítulos y secciones} \\
\midrule
Ya cubierto en el sílabo de econometría & \CT{5--10} (máxima verosimilitud, GMM, pruebas, especificación, métodos semiparamétricos y optimización numérica), \CT{4.6} (regresión cuantílica); \WD{1--9} y \WD{12--14}; ecuaciones simultáneas y SUR (Bloque 11 de ese sílabo) \\
Variable dependiente limitada básica (consulta, sin horas) & \CT{14}, \CT{15.9--15.10}; \WD{15.1--15.6}, \WD{15.9--15.10}, \WD{16.1--16.7}, \WD{17.1--17.6}, \WD{19.1--19.4}; Train (2009), cap.~8. Ya están en el Bloque 12 del sílabo de econometría \\
Métodos bayesianos y de simulación & \CT{12--13} (consulta en el tema 10.3) \\
Sistemas con componentes de error y PCSE & \BT{6}; Beck y Katz (1995). Ver Anexo D.1 \\
Otros modelos & \CT{18--19} (mezclas y riesgos múltiples); \CT{22.8} (modelos lineales mixtos); \CT{23.8}; \WD{18.5}; \WD{20} (consulta en el tema 11.4) \\
Paneles no estacionarios y espaciales & \BT{12--13}; sesiones 4.3 y 4.4 de 1ECO77. Ver Anexo D.1 \\
Otras secciones de panel & \BT{10.2}, \BT{10.4--10.5}; \AR{5} (estructuras de covarianza dinámicas) y apéndices A y B (GMM y extremos, ya en el sílabo de econometría) \\
Regresión cuantílica y efectos por cuantiles & \MHE{7}. Ver Anexo D.4 \\
Ejemplo extendido & \MM{6} (consulta en el tema 4.1) \\
\bottomrule
\end{tabularx}
}

\medskip
\noindent Dos archivos de la carpeta no se usaron porque son imágenes sin
texto: el libro de Maddala (\emph{Limited-Dependent and Qualitative Variables
in Econometrics}, carpeta Panel Data) y \texttt{Sílabo 2026-1.pdf}.

%------------------------------------------------------------------------------
\section*{Anexo C. Programas de la PUCP utilizados}
\addcontentsline{toc}{section}{Anexo C. Programas de la PUCP utilizados}

\subsection*{C.1 Econometría Intermedia: Micro (1ECO25, 2025-1, J.~M.~Rentería)}
Es el programa que guía el plan.

\medskip
{\small\filasalternas
\noindent\begin{tabularx}{\textwidth}{>{\raggedright\arraybackslash}p{6.2cm}X}
\toprule
\rowcolor{tonoOscuro}\cab{Tema del programa} & \cab{Temas del plan} \\
\midrule
1--2. Modelo de regresión lineal clásico & P.2 y sílabo de econometría \\
3. Modelos binarios y de elección múltiple & Bloque 12 del sílabo de econometría y Bloque 10 \\
4. Inferencia causal en economía & Bloque 1 \\
5. Experimentos aleatorios & Bloque 2 \\
6. Emparejamiento & Bloque 3 \\
7. Efectos fijos & 6.2 \\
8. Estudios de eventos & 8.3 \\
9. Dobles diferencias & Bloque 8 \\
10. Variables instrumentales & Bloque 4 \\
11. Regresión discontinua & Bloque 5 \\
12. Otros modelos para datos de panel & 6.3, 6.4, Bloque 7 \\
13. Variable truncada y censurada & 11.1, 11.2 \\
14. Ecuaciones simultáneas y SUR & Bloque 11 del sílabo de econometría; con componentes de error, Anexo D.1 \\
\bottomrule
\end{tabularx}
}

\subsection*{C.2 Otros programas}
{\small\filasalternas
\noindent\begin{tabularx}{\textwidth}{>{\raggedright\arraybackslash}p{6.2cm}X}
\toprule
\rowcolor{tonoOscuro}\cab{Programa} & \cab{Temas del plan} \\
\midrule
1ECO25 (2025-2, L.~García): ecuaciones simultáneas, variables dependientes limitadas, datos de panel y evaluación de impacto & Bloques 1 a 7, 10 y 11 \\
1ECO77, Corte Transversal y Panel (maestría, 2025-2 y 2026-2) & M1: Bloques 1, 3, 5 y 8; M2: Bloque 4; M3: Bloque 6 y tema 7.1; M4: Bloques 7 y 12; paneles no estacionarios y espaciales: Anexo D.1 \\
ECO315, Tópicos de Microeconomía Avanzada (2022-0, C.~Huaroto) & Bloques 2 a 5, 8 y 9; método de replicación: Anexo D.5 \\
ECO387, Muestras complejas e inferencia causal (2022-2) & 3.2, 13.1 \\
ECO386, Inferencia causal y aprendizaje automático (2024-2) & 1.3, 3.5; Anexo D.3 \\
2ECO01, Unidad 3 (maestría, 2026-1) & Bloques 10 y 11 \\
\bottomrule
\end{tabularx}
}

%------------------------------------------------------------------------------
\section*{Anexo D. Material opcional}
\addcontentsline{toc}{section}{Anexo D. Material opcional}

Todo lo de este anexo es opcional: no entra en el cómputo de horas ni modifica
el cronograma.

\subsection*{D.1 Paneles no estacionarios, heterogéneos y espaciales, y sistemas con componentes de error}
Sigue las sesiones 4.3 y 4.4 de 1ECO77. Momento sugerido: después del Bloque 7
y, para la parte no estacionaria, después del Bloque 4 del sílabo de
macroeconometría.
\begin{itemize}[leftmargin=1.4em,itemsep=1pt]
  \item Raíces unitarias y cointegración en paneles: \BT{12}; Pesaran (2015),
        cap.~31.
  \item Paneles heterogéneos grandes y dependencia transversal: Pesaran (2015),
        caps.~28--29.
  \item Paneles espaciales: \BT{13}; Pesaran (2015), cap.~30; Croissant y
        Millo (2019). Enlaza con la econometría espacial del sílabo de
        econometría (Anexo D.3).
  \item Sistemas de ecuaciones con componentes de error (SUR y ecuaciones
        simultáneas): \BT{6}; \BT{7.1--7.3}. Salieron del Bloque 6.
  \item Datos de corte transversal de series de tiempo (pocas unidades
        observadas durante muchos períodos): Beck y Katz (1995) y los errores
        estándar corregidos para panel (PCSE), la alternativa al MCG factible
        cuando $T$ es grande.
\end{itemize}

\subsection*{D.2 Estimación estructural}
Adams (2021), parte II: demanda, selección, demanda con variables
instrumentales, juegos y subastas; \CT{16.7--16.8} (modelo de Roy y modelos
estructurales); Train (2009), caps.~11--14.

\subsection*{D.3 Aprendizaje automático causal}
Sigue el curso ECO386. Ver el Anexo D.1 del sílabo de econometría, que
desarrolla este módulo con Huber, Hansen y los artículos del programa.

\subsection*{D.4 Regresión cuantílica y efectos de tratamiento por cuantiles}
\MHE{7}; \CT{4.6}; tema 13.5 del sílabo de econometría. Complementa las regresiones RIF del tema 3.4.

\subsection*{D.5 Taller de replicación}
Sigue el método del curso ECO315: un artículo por diseño, con los datos de los
autores.
\begin{enumerate}[leftmargin=1.6em,itemsep=2pt]
  \item \textbf{Experimento}: Dammert, Galdo y Galdo (2014).
  \item \textbf{Variables instrumentales}: Angrist y Krueger (1991), con los
        datos de \emph{Mostly Harmless Econometrics}.
  \item \textbf{Regresión discontinua}: los ejemplos de Cattaneo, Idrobo y
        Titiunik (2020).
  \item \textbf{Diferencias en diferencias}: Card y Krueger (1994) y Galdo (2013).
  \item \textbf{Datos de panel}: los ejemplos de Croissant y Millo (2019).
\end{enumerate}
Herramientas: los códigos de \emph{The Remix} (R, Stata y Python), Huber (R),
Adams (2021) (R) y la guía de lecturas de Cai (2024), que reúne la literatura
reciente sobre cada método (variables instrumentales con exposición a
choques, aglomeración, estadísticos suficientes y descomposiciones).

\subsection*{D.6 Software por método}
Paquetes de uso extendido para cada diseño. En Stata se instalan con
\texttt{ssc install} o desde el sitio de sus autores; en R, desde CRAN o
GitHub. La referencia en Stata es Cameron y Trivedi (2022); en Python,
\texttt{pyfixest}, \texttt{DoubleML} y \texttt{EconML} cubren buena parte de
la tabla.

\medskip
{\small\filasalternas
\noindent\begin{tabularx}{\textwidth}{>{\raggedright\arraybackslash}p{4.2cm}>{\raggedright\arraybackslash}X>{\raggedright\arraybackslash}X}
\toprule
\rowcolor{tonoOscuro}\cab{Método (temas)} & \cab{Stata} & \cab{R} \\
\midrule
Efectos fijos de alta dimensión (todo el plan) & \texttt{reghdfe} & \texttt{fixest} \\
Inferencia por aleatorización y pruebas múltiples (2.2) & \texttt{ritest}, \texttt{rwolf}, \texttt{wyoung} & \texttt{ri2}, \texttt{p.adjust} \\
Sensibilidad (3.3) & \texttt{psacalc}, \texttt{sensemakr} & \texttt{sensemakr} \\
Descomposiciones (3.4) & \texttt{oaxaca}, \texttt{rifhdreg} & \texttt{oaxaca} \\
Aprendizaje automático causal (3.5) & \texttt{ddml} & \texttt{DoubleML}, \texttt{grf} \\
Variables instrumentales (4.3) & \texttt{ivreg2}, \texttt{weakivtest}, \texttt{weakiv} & \texttt{fixest}, \texttt{ivmodel} \\
\emph{Shift-share} (4.4) & \texttt{ssaggregate}, \texttt{bartik\_weight} & \texttt{ShiftShareSE}, \texttt{bartik.weight} \\
Regresión discontinua (Bloque 5) & \texttt{rdrobust}, \texttt{rddensity}, \texttt{rdlocrand} & \texttt{rdrobust}, \texttt{rddensity}, \texttt{rdlocrand} \\
Paneles dinámicos (Bloque 7) & \texttt{xtabond2} & \texttt{plm} (\texttt{pgmm}) \\
Diferencias en diferencias (Bloque 8) & \texttt{csdid}, \texttt{did\_imputation}, \texttt{eventstudyinteract}, \texttt{did\_multiplegt}, \texttt{drdid}, \texttt{honestdid} & \texttt{did}, \texttt{didimputation}, \texttt{fixest} (\texttt{sunab}), \texttt{DIDmultiplegt}, \texttt{DRDID}, \texttt{HonestDiD} \\
Control sintético (Bloque 9) & \texttt{synth}, \texttt{sdid} & \texttt{Synth}, \texttt{synthdid} \\
Inferencia con pocos conglomerados (8.2, 13.1) & \texttt{boottest} & \texttt{fwildclusterboot} \\
\bottomrule
\end{tabularx}
}

%------------------------------------------------------------------------------
\section*{Anexo E. Referencias}
\addcontentsline{toc}{section}{Anexo E. Referencias}

\subsection*{E.1 Textos del plan}
{\small
\begin{itemize}[leftmargin=1.4em,itemsep=1pt,label={}]
\item Angrist, J. D. y J.-S. Pischke (2009). \emph{Mostly Harmless Econometrics: An Empiricist's Companion}. Princeton University Press.
\item Angrist, J. D. y J.-S. Pischke (2015). \emph{Mastering 'Metrics: The Path from Cause to Effect}. Princeton University Press.
\item Arellano, M. (2003). \emph{Panel Data Econometrics}. Oxford University Press.
\item Baltagi, B. H. (2021). \emph{Econometric Analysis of Panel Data}, 6.ª ed. Springer.
\item Cameron, A. C. y P. K. Trivedi (2005). \emph{Microeconometrics: Methods and Applications}. Cambridge University Press.
\item Cunningham, S. \emph{Causal Inference: The Remix}, 2.ª ed., en línea: \url{https://mixtape.scunning.com}. Primera edición: \emph{Causal Inference: The Mixtape}, Yale University Press, 2021.
\item Huber, M. (2023). \emph{Causal Analysis: Impact Evaluation and Causal Machine Learning with Applications in R}. MIT Press.
\item Wooldridge, J. M. (2002). \emph{Econometric Analysis of Cross Section and Panel Data}. MIT Press.
\end{itemize}
}

\subsection*{E.2 Libros y textos de consulta de la carpeta (sin horas, salvo donde se indica)}
{\small
\begin{itemize}[leftmargin=1.4em,itemsep=1pt,label={}]
\item Adams, C. P. (2021). \emph{Learning Microeconometrics with R}. CRC Press.
\item Bailey, M. A. (2020). \emph{Real Econometrics: The Right Tools to Answer Important Questions}, 2.ª ed. Oxford University Press.
\item Beltrán, A. y J. F. Castro (2010). \emph{Modelos de datos de panel y variables dependientes limitadas: teoría y práctica}. Universidad del Pacífico.
\item Biørn, E. (2017). \emph{Econometrics of Panel Data: Methods and Applications}. Oxford University Press.
\item Cai, C. (2024). Literature on Recent Advances in Applied Micro Methods. Guía de lecturas, versión del 9 de junio de 2024.
\item Croissant, Y. y G. Millo (2019). \emph{Panel Data Econometrics with R}. Wiley.
\item García Núñez, L. (2010). Econometría de evaluación de impacto. Documento de Trabajo 283, Departamento de Economía, PUCP.
\item Gertler, P., S. Martínez, P. Premand, L. Rawlings y C. Vermeersch (2017). \emph{La evaluación de impacto en la práctica}, 2.ª ed. Banco Mundial y BID. Entra en las horas en el tema 2.3.
\item Hansen, B. E. (2021). \emph{Econometrics}. Manuscrito, University of Wisconsin (carpeta de econometría).
\item Hsiao, C. (2003). \emph{Analysis of Panel Data}, 2.ª ed. Cambridge University Press.
\item Neal, B. (2020). \emph{Introduction to Causal Inference from a Machine Learning Perspective}. Notas de curso.
\item Pesaran, M. H. (2015). \emph{Time Series and Panel Data Econometrics}. Oxford University Press.
\item Train, K. E. (2009). \emph{Discrete Choice Methods with Simulation}, 2.ª ed. Cambridge University Press. Entra en las horas en los temas 10.1 y 10.3.
\item Wooldridge, J. M. (2012). \emph{Introductory Econometrics: A Modern Approach}, 5.ª ed. Cengage (carpeta de econometría).
\end{itemize}
}

\subsection*{E.3 Referencias tomadas de los programas de la PUCP}
{\small
\begin{itemize}[leftmargin=1.4em,itemsep=1pt,label={}]
\item Abadie, A., S. Athey, G. Imbens y J. Wooldridge (2017). When Should You Adjust Standard Errors for Clustering? NBER Working Paper 24003.
\item Abadie, A. y G. Imbens (2016). Matching on the Estimated Propensity Score. \emph{Econometrica}, 84(2), 781--807.
\item Anderson, T. W. y H. Rubin (1949). Estimation of the Parameters of a Single Equation in a Complete System of Stochastic Equations. \emph{Annals of Mathematical Statistics}, 20(1), 46--63.
\item Angrist, J. D. (2022). Empirical Strategies in Economics: Illuminating the Path from Cause to Effect. \emph{Econometrica}, 90(6), 2509--2539.
\item Angrist, J. D. y A. B. Krueger (1991). Does Compulsory School Attendance Affect Schooling and Earnings? \emph{Quarterly Journal of Economics}, 106(4), 979--1014.
\item Bédécarrats, F., I. Guérin y F. Roubaud (eds.) (2020). \emph{Randomized Control Trials in the Field of Development: A Critical Perspective}. Oxford University Press.
\item Caliendo, M. y S. Kopeinig (2008). Some Practical Guidance for the Implementation of Propensity Score Matching. \emph{Journal of Economic Surveys}, 22, 31--72.
\item Callaway, B. y P. H. C. Sant'Anna (2021). Difference-in-Differences with Multiple Time Periods. \emph{Journal of Econometrics}, 225(2), 200--230.
\item Cameron, A. C. y D. L. Miller (2015). A Practitioner's Guide to Cluster-Robust Inference. \emph{Journal of Human Resources}, 50(2), 317--372.
\item Card, D. y A. B. Krueger (1994). Minimum Wages and Employment: A Case Study of the Fast-Food Industry in New Jersey and Pennsylvania. \emph{American Economic Review}, 84(4), 772--793.
\item Cattaneo, M. D., N. Idrobo y R. Titiunik (2020). \emph{A Practical Introduction to Regression Discontinuity Designs: Foundations}. Cambridge University Press.
\item Cattaneo, M. D. y R. Titiunik (2022). Regression Discontinuity Designs. \emph{Annual Review of Economics}, 14, 821--851.
\item Chay, K. Y. y M. Greenstone (2005). Does Air Quality Matter? Evidence from the Housing Market. \emph{Journal of Political Economy}, 113(2), 376--424.
\item Chernozhukov, V., D. Chetverikov, M. Demirer, E. Duflo, C. Hansen, W. Newey y J. Robins (2018). Double/Debiased Machine Learning for Treatment and Structural Parameters. \emph{The Econometrics Journal}, 21(1).
\item Cinelli, C., A. Forney y J. Pearl (2022). A Crash Course in Good and Bad Controls. \emph{Sociological Methods \& Research}.
\item Dammert, A., J. Galdo y V. Galdo (2014). Preventing Dengue through Mobile Phones: Evidence from a Field Experiment in Peru. \emph{Journal of Health Economics}, 35, 147--161.
\item de Chaisemartin, C. y X. D'Haultfœuille (2020). Two-Way Fixed Effects Estimators with Heterogeneous Treatment Effects. \emph{American Economic Review}, 110(9), 2964--2996.
\item Dehejia, R. y S. Wahba (2002). Propensity Score-Matching Methods for Nonexperimental Causal Studies. \emph{Review of Economics and Statistics}, 84(1), 151--161.
\item Duflo, E., R. Glennerster y M. Kremer (2006). Using Randomization in Development Economics Research: A Toolkit. MIT Department of Economics Working Paper 06-36.
\item Evans, W. N. y R. M. Schwab (1995). Finishing High School and Starting College: Do Catholic Schools Make a Difference? \emph{Quarterly Journal of Economics}, 110(4), 941--974.
\item Freyaldenhoven, S., C. Hansen, J. Pérez y J. Shapiro (2024). Visualization, Identification, and Estimation in the Linear Panel Event-Study Design. NBER Working Paper 29170.
\item Galdo, J. (2013). Long-Run Labor-Market Impacts of Early-Life Exposure to Civil War: Evidence from Shining Path in Peru. \emph{Economic Development and Cultural Change}, 61(4), 789--823.
\item Glennerster, R. y K. Takavarasha (2013). \emph{Running Randomized Evaluations: A Practical Guide}. Princeton University Press.
\item Goodman-Bacon, A. (2021). Difference-in-Differences with Variation in Treatment Timing. \emph{Journal of Econometrics}, 225(2), 254--277.
\item Heckman, J. y R. Pinto (2024). Econometric Causality: The Central Role of Thought Experiments. \emph{Journal of Econometrics}, 243(1--2), 105719.
\item Holland, P. W. (1986). Statistics and Causal Inference. \emph{Journal of the American Statistical Association}, 81(396), 945--960.
\item Imai, K. e I. S. Kim (2019). When Should We Use Unit Fixed Effects Regression Models for Causal Inference with Longitudinal Data? \emph{American Journal of Political Science}, 63(2), 467--490.
\item Imbens, G. (2020). Potential Outcome and Directed Acyclic Graph Approaches to Causality. \emph{Journal of Economic Literature}, 58(4), 1129--1179.
\item Imbens, G. (2022). Causality in Econometrics: Choice vs Chance. \emph{Econometrica}, 90(6), 2541--2566.
\item Imbens, G. y T. Lemieux (2007). Regression Discontinuity Designs: A Guide to Practice. NBER Working Paper 13039.
\item LaLonde, R. J. (1986). Evaluating the Econometric Evaluations of Training Programs with Experimental Data. \emph{American Economic Review}, 76(4), 604--620.
\item Miller, D. (2023). An Introductory Guide to Event Study Models. \emph{Journal of Economic Perspectives}, 37(2), 203--230.
\item Mogstad, M. y A. Torgovitsky (2024). Instrumental Variables with Unobserved Heterogeneity in Treatment Effects. En C. Dustmann y T. Lemieux (eds.), \emph{Handbook of Labor Economics}, vol.~5, cap.~1.
\item Rosenbaum, P. R. y D. B. Rubin (1983). The Central Role of the Propensity Score in Observational Studies for Causal Effects. \emph{Biometrika}, 70(1), 41--55.
\item Roth, J. (2022). Pretest with Caution: Event-Study Estimates after Testing for Parallel Trends. \emph{American Economic Review: Insights}, 4(3), 305--322.
\item Roth, J., P. H. C. Sant'Anna, A. Bilinski y J. Poe (2023). What's Trending in Difference-in-Differences? A Synthesis of the Recent Econometrics Literature. \emph{Journal of Econometrics}, 235(2), 2218--2244.
\item Rubin, D. B. (1974). Estimating Causal Effects of Treatments in Randomized and Nonrandomized Studies. \emph{Journal of Educational Psychology}, 66(5), 688--701.
\item Sant'Anna, P. H. C. y J. Zhao (2020). Doubly Robust Difference-in-Differences Estimators. \emph{Journal of Econometrics}, 219(1), 101--122.
\item Smith, J. y P. Todd (2005). Does Matching Overcome LaLonde's Critique of Nonexperimental Estimators? \emph{Journal of Econometrics}, 125(1--2), 305--353.
\item Solon, G., S. J. Haider y J. M. Wooldridge (2015). What Are We Weighting For? \emph{Journal of Human Resources}, 50(2), 301--316.
\item Stock, J. H. y M. Yogo (2005). Testing for Weak Instruments in Linear IV Regression. En D. W. K. Andrews y J. H. Stock (eds.), \emph{Identification and Inference for Econometric Models}. Cambridge University Press.
\item Sun, L. y S. Abraham (2021). Estimating Dynamic Treatment Effects in Event Studies with Heterogeneous Treatment Effects. \emph{Journal of Econometrics}, 225(2), 175--199.
\end{itemize}
}

\subsection*{E.4 Otras referencias (incluye las añadidas en la revisión)}
{\small
\begin{itemize}[leftmargin=1.4em,itemsep=1pt,label={}]
\item Abadie, A. (2021). Using Synthetic Controls: Feasibility, Data Requirements, and Methodological Aspects. \emph{Journal of Economic Literature}, 59(2), 391--425.
\item Abadie, A., A. Diamond y J. Hainmueller (2010). Synthetic Control Methods for Comparative Case Studies: Estimating the Effect of California's Tobacco Control Program. \emph{Journal of the American Statistical Association}, 105(490), 493--505.
\item Anderson, M. L. (2008). Multiple Inference and Gender Differences in the Effects of Early Intervention: A Reevaluation of the Abecedarian, Perry Preschool, and Early Training Projects. \emph{Journal of the American Statistical Association}, 103(484), 1481--1495.
\item Andrews, I., J. H. Stock y L. Sun (2019). Weak Instruments in Instrumental Variables Regression: Theory and Practice. \emph{Annual Review of Economics}, 11, 727--753.
\item Arkhangelsky, D., S. Athey, D. A. Hirshberg, G. W. Imbens y S. Wager (2021). Synthetic Difference-in-Differences. \emph{American Economic Review}, 111(12), 4088--4118.
\item Athey, S. y G. W. Imbens (2017). The Econometrics of Randomized Experiments. En A. V. Banerjee y E. Duflo (eds.), \emph{Handbook of Economic Field Experiments}, vol.~1, 73--140. Elsevier.
\item Beck, N. y J. N. Katz (1995). What to Do (and Not to Do) with Time-Series Cross-Section Data. \emph{American Political Science Review}, 89(3), 634--647.
\item Blinder, A. S. (1973). Wage Discrimination: Reduced Form and Structural Estimates. \emph{Journal of Human Resources}, 8(4), 436--455.
\item Borusyak, K., P. Hull y X. Jaravel (2022). Quasi-Experimental Shift-Share Research Designs. \emph{Review of Economic Studies}, 89(1), 181--213.
\item Borusyak, K., X. Jaravel y J. Spiess (2024). Revisiting Event-Study Designs: Robust and Efficient Estimation. \emph{Review of Economic Studies}, 91(6), 3253--3285.
\item Bruhn, M. y D. McKenzie (2009). In Pursuit of Balance: Randomization in Practice in Development Field Experiments. \emph{American Economic Journal: Applied Economics}, 1(4), 200--232.
\item Cameron, A. C. y P. K. Trivedi (2022). \emph{Microeconometrics Using Stata}, 2.ª ed., 2 vols. Stata Press.
\item Cattaneo, M. D., N. Idrobo y R. Titiunik (2024). \emph{A Practical Introduction to Regression Discontinuity Designs: Extensions}. Cambridge University Press.
\item Cinelli, C. y C. Hazlett (2020). Making Sense of Sensitivity: Extending Omitted Variable Bias. \emph{Journal of the Royal Statistical Society, Series B}, 82(1), 39--67.
\item DiNardo, J., N. M. Fortin y T. Lemieux (1996). Labor Market Institutions and the Distribution of Wages, 1973--1992: A Semiparametric Approach. \emph{Econometrica}, 64(5), 1001--1044.
\item Firpo, S., N. M. Fortin y T. Lemieux (2009). Unconditional Quantile Regressions. \emph{Econometrica}, 77(3), 953--973.
\item Fortin, N., T. Lemieux y S. Firpo (2011). Decomposition Methods in Economics. En O. Ashenfelter y D. Card (eds.), \emph{Handbook of Labor Economics}, vol.~4A, cap.~1, 1--102. Elsevier.
\item Gelman, A. y G. Imbens (2019). Why High-Order Polynomials Should Not Be Used in Regression Discontinuity Designs. \emph{Journal of Business \& Economic Statistics}, 37(3), 447--456.
\item Goldsmith-Pinkham, P., I. Sorkin y H. Swift (2020). Bartik Instruments: What, When, Why, and How. \emph{American Economic Review}, 110(8), 2586--2624.
\item Lee, D. S., J. McCrary, M. J. Moreira y J. Porter (2022). Valid $t$-Ratio Inference for IV. \emph{American Economic Review}, 112(10), 3260--3290.
\item List, J. A., A. M. Shaikh e Y. Xu (2019). Multiple Hypothesis Testing in Experimental Economics. \emph{Experimental Economics}, 22, 773--793.
\item Montiel Olea, J. L. y C. Pflueger (2013). A Robust Test for Weak Instruments. \emph{Journal of Business \& Economic Statistics}, 31(3), 358--369.
\item Oaxaca, R. (1973). Male-Female Wage Differentials in Urban Labor Markets. \emph{International Economic Review}, 14(3), 693--709.
\item Oster, E. (2019). Unobservable Selection and Coefficient Stability: Theory and Evidence. \emph{Journal of Business \& Economic Statistics}, 37(2), 187--204.
\item Rambachan, A. y J. Roth (2023). A More Credible Approach to Parallel Trends. \emph{Review of Economic Studies}, 90(5), 2555--2591.
\item Roodman, D. (2009a). How to Do xtabond2: An Introduction to Difference and System GMM in Stata. \emph{The Stata Journal}, 9(1), 86--136.
\item Roodman, D. (2009b). A Note on the Theme of Too Many Instruments. \emph{Oxford Bulletin of Economics and Statistics}, 71(1), 135--158.
\item Young, A. (2019). Channeling Fisher: Randomization Tests and the Statistical Insignificance of Seemingly Significant Experimental Results. \emph{Quarterly Journal of Economics}, 134(2), 557--598.
\end{itemize}
}

\vspace{1.2cm}
\franjas[2pt]
\begin{center}\small\color{tonoPizarra}
Plan de lectura de elaboración propia. Los títulos de los capítulos
corresponden a las obras citadas y se reproducen únicamente con fines de
referencia bibliográfica. La asignación de capítulos y secciones se basa en
los índices de los libros de la carpeta y, para \emph{The Remix}, en el índice
de su sitio web; los temas de la PUCP proceden de los programas citados en el
Anexo C.
\end{center}

\end{document}
